\documentclass[10pt, xcolor={dvipsnames}]{beamer}

%========================
% VIETNAMESE & PACKAGES
%========================
\usepackage[utf8]{inputenc}
\usepackage[T5]{fontenc}
\usepackage[vietnamese]{babel}
\usepackage{tikz}
\usetikzlibrary{arrows.meta, backgrounds, positioning, shadows.blur, shapes.geometric, calc}
\usepackage{listings}
\usepackage{amsmath, amssymb}
\usepackage{tcolorbox}
\tcbuselibrary{skins, shadings}

%========================
% THEME & VIBRANT COLORS
%========================
\usetheme{Madrid}
\definecolor{mainblue}{RGB}{0, 102, 204}
\definecolor{darkblue}{RGB}{0, 32, 96}
\definecolor{lightorange}{RGB}{255, 165, 0}
\definecolor{softblue}{RGB}{235, 245, 255}

\setbeamercolor{palette primary}{bg=mainblue, fg=white}
\setbeamercolor{palette secondary}{bg=lightorange, fg=white}
\setbeamercolor{palette tertiary}{bg=white, fg=mainblue}
\setbeamercolor{structure}{fg=mainblue}

% Khung nội dung đẹp
\newtcolorbox{beautybox}[2]{
    enhanced, colback=white, colframe=#1, fonttitle=\bfseries, title=#2,
    boxrule=1pt, arc=10pt, shadow={1pt}{-1pt}{0pt}{black!10},
    attach boxed title to top left={xshift=5mm, yshift=-2mm},
    boxed title style={colback=#1}
}

% Style Code C++
\lstset{
    language=C++,
    basicstyle=\ttfamily\scriptsize,
    keywordstyle=\color{blue}\bfseries,
    commentstyle=\color{gray},
    stringstyle=\color{orange},
    backgroundcolor=\color{softblue},
    frame=none, breaklines=true
}

\begin{document}

% 1. Slide Tiêu đề - Hiện đại & Chuyên nghiệp
\begin{frame}[plain]
    \begin{tikzpicture}[remember picture, overlay]
        \shade[top color=darkblue, bottom color=mainblue!60!black] (current page.south west) rectangle (current page.north east);
        \fill[lightorange, opacity=0.15] (current page.north east) -- (current page.north west) -- ($(current page.north west) + (0,-3)$) -- cycle;
        \node[fill=white, rounded corners=15pt, inner sep=20pt, shadow={3pt}{-3pt}{10pt}{black!40}] at (current page.center) {
            \begin{minipage}{0.85\textwidth}
                \centering
                \Huge \textcolor{darkblue}{\textbf{LÝ THUYẾT ĐỒ THỊ}}\\
                \vspace{0.4cm}
                \large \textcolor{lightorange!80!black}{\textbf{CHIẾN THUẬT QUÉT SẠCH T-MATH}}\\
                \vspace{0.8cm}
                \small \textcolor{gray}{Giảng Viên: \textbf{Đặng Xuân Lâm}}
            \end{minipage}
        };
    \end{tikzpicture}
\end{frame}

% 2. Mục lục
\begin{frame}{Hành trình chinh phục}
    \tableofcontents
\end{frame}

\section{Tổng quan Đồ thị}
% 3. Định nghĩa
\begin{frame}{Đồ thị là gì?}
    \begin{beautybox}{mainblue}{Định nghĩa cốt lõi}
        Đồ thị $G = (V, E)$ gồm tập các đỉnh (Vertices) và tập các cạnh (Edges) kết nối các cặp đỉnh.
    \end{beautybox}
    \begin{columns}
        \begin{column}{0.5\textwidth}
            \begin{itemize}
                \item Đỉnh: Tài khoản Facebook.
                \item Cạnh: Kết bạn, nhắn tin.
            \end{itemize}
        \end{column}
        \begin{column}{0.5\textwidth}
            \centering
            \begin{tikzpicture}[scale=0.8]
                \foreach \p/\n in {(0,0)/1, (2,1)/2, (2,-1)/3, (4,0)/4}
                    \node[circle, draw=mainblue, fill=softblue, thick] (\n) at \p {\n};
                \draw[thick] (1)--(2); \draw[thick] (1)--(3); \draw[thick] (2)--(4); \draw[thick] (3)--(4);
            \end{tikzpicture}
        \end{column}
    \end{columns}
\end{frame}

% 4. Phân loại
\begin{frame}{Phân loại: Có hướng \& Vô hướng}
    \begin{columns}
        \begin{column}{0.5\textwidth}
            \begin{beautybox}{lightorange}{Vô hướng}
                Đường 2 chiều. Bạn bè trên FB là vô hướng.
            \end{beautybox}
            \centering
            \begin{tikzpicture}\node[circle,draw](A){A};\node[circle,draw,right of=A](B){B};\draw(A)--(B);\end{tikzpicture}
        \end{column}
        \begin{column}{0.5\textwidth}
            \begin{beautybox}{mainblue}{Có hướng}
                Đường 1 chiều. Follow trên TikTok là có hướng.
            \end{beautybox}
            \centering
            \begin{tikzpicture}\node[circle,draw](A){A};\node[circle,draw,right of=A](B){B};\draw[-{Stealth}](A)--(B);\end{tikzpicture}
        \end{column}
    \end{columns}
\end{frame}

\section{Lưu trữ (Memory Optimization)}
% 5. Danh sách kề
\begin{frame}[fragile]{Lưu trữ: Danh sách kề (Adjacency List)}
    \begin{beautybox}{mainblue}{Vũ khí tối ưu}
        Tiết kiệm bộ nhớ khi đồ thị có $10^5$ đỉnh. Duyệt cực nhanh!
    \end{beautybox}
    \begin{lstlisting}
vector<int> adj[100005];
void add_edge(int u, int v) {
    adj[u].push_back(v);
    adj[v].push_back(u); 
}
    \end{lstlisting}
\end{frame}

% 6. Ma trận kề
\begin{frame}{Lưu trữ: Ma trận kề (Adjacency Matrix)}
    \begin{itemize}
        \item Dùng mảng 2 chiều \texttt{A[N][N]}.
        \item \textbf{Ưu điểm:} Kiểm tra $u, v$ có kề nhau không trong $O(1)$.
        \item \textbf{Nhược điểm:} Tốn bộ nhớ $O(N^2)$, dễ bị \textbf{MLE} nếu $N > 10^4$.
    \end{itemize}
\end{frame}

\section{Thuật toán Tìm kiếm}
% 7. DFS
\begin{frame}{DFS: Đi xuyên bóng tối (Depth First Search)}
    \begin{beautybox}{mainblue}{Chiến thuật Đệ quy}
        Ưu tiên đi sâu nhất có thể. Dùng \textbf{Stack} hoặc \textbf{Đệ quy}.
    \end{beautybox}
    \begin{itemize}
        \item Ứng dụng: Tìm thành phần liên thông, kiểm tra chu trình.
        \item Độ phức tạp: $O(V + E)$.
    \end{itemize}
\end{frame}

% 8. BFS
\begin{frame}{BFS: Lan tỏa sức mạnh (Breadth First Search)}
    \begin{beautybox}{lightorange}{Chiến thuật Hàng đợi}
        Quét theo từng lớp (loang). Dùng \textbf{Queue}.
    \end{beautybox}
    \begin{itemize}
        \item Ứng dụng: Tìm \textbf{đường đi ngắn nhất} trên đồ thị không trọng số.
        \item Độ phức tạp: $O(V + E)$.
    \end{itemize}
\end{frame}

\section{Bài tập & Cách giải (T-Math Style)}
% 9. Thành phần liên thông
\begin{frame}{Bài toán 1: Đếm số vùng liên thông}
    \begin{itemize}
        \item \textbf{Bài toán:} Đếm số "đảo" trong đồ thị.
        \item \textbf{Cách giải:} 
            1. Duyệt tất cả các đỉnh.
            2. Nếu đỉnh $i$ chưa thăm, tăng biến đếm và gọi DFS/BFS($i$) để đánh dấu toàn bộ đảo đó.
    \end{itemize}
\end{frame}

% 10. Mê cung 2D
\begin{frame}{Bài toán 2: Mê cung 2D}
    \begin{beautybox}{lightorange}{Mẹo xử lý}
        Dùng mảng hướng để code cực nhanh:
        \texttt{dx[] = \{-1, 1, 0, 0\}}, \texttt{dy[] = \{0, 0, -1, 1\}}.
    \end{beautybox}
    \centering
    \begin{tikzpicture}[scale=0.5]
        \draw[step=1,gray,very thin] (0,0) grid (4,3);
        \node at (0.5,2.5) {S}; \node at (3.5,0.5) {E};
        \fill[red!20] (1,1) rectangle (2,2); \fill[red!20] (2,2) rectangle (3,3);
    \end{tikzpicture}
\end{frame}

\section{Thuật toán Nâng cao}
% 11. Dijkstra
\begin{frame}{Dijkstra: Đường đi ngắn nhất}
    \begin{beautybox}{mainblue}{Tham lam (Greedy)}
        Luôn chọn đỉnh có khoảng cách nhỏ nhất để tối ưu tiếp.
    \end{beautybox}
    \textbf{Độ phức tạp:} $O(E \log V)$ khi dùng Priority Queue.
\end{frame}

% 12. Kruskal
\begin{frame}{Kruskal: Cây khung nhỏ nhất (MST)}
    \begin{itemize}
        \item Sắp xếp các cạnh theo trọng số tăng dần.
        \item Dùng \textbf{DSU (Disjoint Set Union)} để kiểm tra chu trình.
        \item Gộp các đỉnh lại cho đến khi đủ $V-1$ cạnh.
    \end{itemize}
\end{frame}

% 13. Fast I/O
\begin{frame}[fragile]{Mẹo đi thi: Tăng tốc code}
    Để tránh bị \textbf{TLE} vô lý trên hệ thống T-Math:
\begin{lstlisting}
ios_base::sync_with_stdio(0);
cin.tie(0); cout.tie(0);
// Dung '\n' thay vi endl
\end{lstlisting}
\end{frame}

\section{ĐẤU TRƯỜNG TRÍ TUỆ (GAME)}
% 14. Game 1
\begin{frame}{GAME: PHÁ ĐẢO (Câu 1/5)}
    \begin{beautybox}{mainblue}{Thử thách}
        Để tìm đường ngắn nhất trên đồ thị \textbf{không trọng số}, ta dùng thuật toán nào?
    \end{beautybox}
    \pause
    \begin{tcolorbox}[colback=green!10, colframe=green!50!black]
        \textbf{Đáp án:} BFS (Hàng đợi).
    \end{tcolorbox}
\end{frame}

% 15. Game 2
\begin{frame}{GAME: PHÁ ĐẢO (Câu 2/5)}
    \begin{beautybox}{lightorange}{Thử thách}
        Đồ thị $10^5$ đỉnh, dùng ma trận kề sẽ bị lỗi gì?
    \end{beautybox}
    \pause
    \begin{tcolorbox}[colback=red!10, colframe=red!50!black]
        \textbf{Đáp án:} MLE (Memory Limit Exceeded - Tràn bộ nhớ).
    \end{tcolorbox}
\end{frame}

% 16. Game 3
\begin{frame}{GAME: PHÁ ĐẢO (Câu 3/5)}
    \begin{beautybox}{mainblue}{Thử thách}
        Thuật toán DFS thường được cài đặt bằng cách sử dụng?
    \end{beautybox}
    \pause
    \begin{tcolorbox}[colback=green!10, colframe=green!50!black]
        \textbf{Đáp án:} Ngăn xếp (Stack) hoặc Đệ quy.
    \end{tcolorbox}
\end{frame}

% 17. Game 4
\begin{frame}{GAME: PHÁ ĐẢO (Câu 4/5)}
    \begin{beautybox}{lightorange}{Thử thách}
        Kruskal dùng để tìm cái gì?
    \end{beautybox}
    \pause
    \begin{tcolorbox}[colback=green!10, colframe=green!50!black]
        \textbf{Đáp án:} Cây khung nhỏ nhất (Minimum Spanning Tree).
    \end{tcolorbox}
\end{frame}

% 18. Game 5
\begin{frame}{GAME: PHÁ ĐẢO (Câu 5 - BOSS)}
    \begin{beautybox}{mainblue}{Thử thách}
        Độ phức tạp tối ưu của Dijkstra là bao nhiêu?
    \end{beautybox}
    \pause
    \begin{tcolorbox}[colback=yellow!10, colframe=yellow!50!black]
        \textbf{Đáp án:} $O(E \log V)$.
    \end{tcolorbox}
\end{frame}

% 19. Tổng kết
\begin{frame}{Bí kíp luyện rồng}
    \begin{itemize}
        \item Luôn vẽ đồ thị ra giấy trước khi gõ code.
        \item Thành thạo DFS/BFS là nắm chắc 50% điểm số.
        \item Luyện tập tại T-Math để "out trình" đối thủ!
    \end{itemize}
\end{frame}

% 20. Lời kết
\begin{frame}[plain]
    \centering
    \Huge \textcolor{mainblue}{\textbf{CHÚC MAY MẮN!}}\\
    \vspace{0.5cm}
    \large \textcolor{gray}{Hẹn gặp lại các bạn ở đỉnh vinh quang T-Math!}
\end{frame}

\end{document}